% !TeX program = xelatex
\documentclass[UTF8,fontset=fandol,12pt]{ctexart}
\usepackage[a4paper,margin=22mm,headheight=16pt]{geometry}
\usepackage{booktabs,tabularx,array,xcolor,fancyhdr,hyperref}
\definecolor{accent}{HTML}{28649A}
\hypersetup{colorlinks=true,linkcolor=accent,urlcolor=accent,pdftitle={第8队研讨汇报：8分钟逐页讲稿},pdfauthor={NLP 第8队}}
% 十二页共用讲稿，供 Beamer notes 与独立讲稿引用。
\newcommand{\speechOne}{老师好，我们是第八队。课设题目是长对话结构化记忆复现与优化。今天主要介绍六篇论文怎样支撑同一研究问题，以及小组协作完成的实验和后续计划。}

\newcommand{\speechTwo}{我们研究长对话问答：根据一个人长期积累的历史，回答经历、关系和时间问题。历史很长，相关信息往往分散。同一个人多次谈旅行，检索可能找到主题相关、却属于另一次旅行的内容。“去年”又必须结合记录年份才能解释。因此，我们关注如何保留事件联系，让模型正确使用证据，同时检查结构化带来的额外成本。}

\newcommand{\speechSynthesis}{六篇论文被放进同一个分析框架：信息怎样组织、证据怎样找到、取回后怎样使用，以及更新后怎样验证。StructMem 用事件绑定和跨事件综合保留关系，是我们的实现主线。Hindsight 强调时间信息和多路检索，Semantic XPath 强调结构范围，帮助我们检查是否找对事件。MetaMem 进一步提醒，证据取到了也可能用错。HyMEM 关注经验更新，Text2Mem 关注操作及执行状态，两者帮助讨论怎样安全地维护记忆，但动态更新暂不纳入本轮实现。由此，其他论文参与定义检查维度和研究边界，而不是只在背景里列名字。它们任务和指标不同，我们也不直接拿论文分数排名。}

\newcommand{\speechThree}{具体实现分三步。先抽取事实和关系，保留说话人、来源消息和时间绑定。然后做跨事件综合：汇集一个时间窗口的新事件，检索历史种子并补齐上下文，把联系写成独立综合记忆。最后根据问题检索事件和综合条目，再交给回答模型。本轮使用一小时窗口、十五个历史种子，每题六十条事件和五条综合。事件绑定和综合来自原论文，我们完成的是适配、运行和分析。}

\newcommand{\speechFour}{实验使用 LoCoMo，固定三个开发会话、三十道题，覆盖表中的四类问题；另外七个会话留待后续评价。参考答案只在生成回答后用于评分。S 零检索六十条原始对话块；S 一只保留事件记忆；S 二再加入跨事件综合。当前阶段比较 S 零和 S 二，S 一用于后续的综合模块消融。两组公共生成参数一致，但记忆表示、输入长度和方法指令不同，因此分差描述系统对照，不能全部归因于某一个模块。}

\newcommand{\speechFive}{当前链路已经从数据贯通到评分：近两千条消息生成三千七百一十二条事件记忆，再形成九十一条综合记忆。S 零和 S 二各完成三十条真实回答，六十条裁判记录的输入指纹与分数已经离线核对。模型替换、关闭压缩等复现差异也保留记录。}

\newcommand{\speechSix}{应用既有证据裁决后，S 零答对十五题，正确率百分之五十；S 二答对十九题，正确率百分之六十三点三。多出的四题中，时间题占两题，开放域和单跳各一题，多跳相同。同时，回答输入从十三万两千多 tokens 增加到五十八万六千多，约四点四二倍，尚未计入记忆构建。原始自动裁判给 S 零十六题正确，其中一条日期回答存在矛盾，经证据复核改为错误，所以正文显示十五题。原始判定和后续裁决分别保留。这些结果只描述固定三十道开发题，还不能证明稳定提升。}

\newcommand{\speechSeven}{用同一个旅行案例，可以把几篇论文的视角连起来。消息记录于二零二三年四月十日，内容说的是“我们去年进行的公路旅行”，相关对话确认了地点。S 零把记录日期当成旅行日期；S 二回答去年，也就是二零二二年。StructMem 关注人物、地点和事件的联系；Hindsight 提醒区分记录时间与发生时间；MetaMem 提醒检查取回的证据是否用对。这里是在用多篇论文共同分析一个真实结果，并不意味着已经实现它们的全部机制，也不能仅凭这一个案例断定改进原因。}

\newcommand{\speechEight}{文献比较与现有结果共同指向下一步：先弄清收益来自哪里，再控制输入成本。首先开展 S 一消融，在事件库、模型和方法指令一致的条件下，只关闭跨事件综合。其次把事件检索数从六十减到三十，综合记忆仍取五条，观察质量和输入量如何变化。如果下降也如实保留。开发集确定配置以后，再从七个留出会话分层抽取约一百题评价。动态更新和基于反馈的规则学习暂不扩展，先把当前问题做扎实。以上都是待运行方案。}

\newcommand{\speechNine}{六人分工围绕同一条链路衔接。A 冻结数据、划分和题单；B 提供普通检索；C 交付事件记忆；D 接着生成综合记忆并完成 S 二；E 评分和分析；F 提供统一回答接口，完成 S 零、联调与复算。连接这些工作的是同一题单、统一输入输出和逐题来源记录。这里不仅有任务划分，也有下游接收与交叉核查。表中列出的是目前大致贡献。}

\newcommand{\speechCollaboration}{我们主要通过 GitHub、微信群和线下讨论协作。GitHub 用来管理任务、代码和文档，通过 Issue 和 PR 记录进展、交接和修改；微信群用于日常沟通、同步进度，及时协调接口、数据和材料问题；线下则集中讨论论文之间的联系、实验方案、结果分析和汇报内容。六位同学虽然负责不同模块，但始终围绕同一个研究问题、同一批题目和同一条实验链路推进，并在交接时相互核对。这样，个人阅读和实现能够汇入共同的分析与实验。}

\newcommand{\speechTen}{最后，我们用共同问题组织六篇论文，再围绕同一实验链路分工、交接和核查。现有结果是 S 二多答对四题，回答输入约四点四二倍。下一步开展综合模块消融和预算对照，并通过留出评价，形成可复跑、有证据的结论。}

\pagestyle{fancy}\fancyhf{}
\fancyhead[L]{\small 第 8 队 · 第一次课设研讨}
\fancyhead[R]{\small 逐页讲稿}
\fancyfoot[C]{\thepage}
\renewcommand{\headrulewidth}{0.3pt}
\setlength{\parindent}{0pt}
\setlength{\parskip}{5pt}
\linespread{1.3}
\newcommand{\talk}[4]{%
  {\color{accent}\large\bfseries 第 #1 页\quad #2}\par
  {\small #3}\par\vspace{2mm}
  #4\par}
\begin{document}
{\LARGE\bfseries 长对话结构化记忆复现与优化}\par
\vspace{2mm}{\Large 8 分钟逐页讲稿}\par
2026 年 10 月 22 日研讨\quad 内容更新至 2026 年 10 月 21 日\par
线下汇报：A（ngocnamd93-spec）\quad 视频录制：F（hycx233）

\vspace{4mm}
\textbf{使用方式}

讲稿对应幻灯片正文第 1--12 页，可直接朗读。补充说明与问答准备见独立的《补充说明与答辩准备》。下表为计划节奏，合计 480 秒；它不是实际录音计时。建议完整排练一次，为翻页、停顿和较慢的数字读法留出时间。

\vspace{4mm}
\begin{tabularx}{\linewidth}{@{}rXrr@{}}
\toprule 页 & 主题 & 本页时长 & 累计时间\\\midrule
1 & 题目与汇报范围 & 15 秒 & 0:15\\
2 & 研究问题 & 35 秒 & 0:50\\
3 & 六篇论文的系统性提炼 & 60 秒 & 1:50\\
4 & StructMem 方法 & 40 秒 & 2:30\\
5 & 实验设计 & 40 秒 & 3:10\\
6 & 实际进展 & 30 秒 & 3:40\\
7 & 初步结果与输入成本 & 55 秒 & 4:35\\
8 & 共同案例与论文视角 & 45 秒 & 5:20\\
9 & 创新性论证与后续方案 & 45 秒 & 6:05\\
10 & 六人分工与贡献 & 40 秒 & 6:45\\
11 & 小组协作方式 & 50 秒 & 7:35\\
12 & 阶段结论与下一步 & 25 秒 & 8:00\\\bottomrule
\end{tabularx}

\vspace{4mm}
\textbf{统一读法与事实口径}
\begin{itemize}
\item S0 / S1 / S2 读作“S 零 / S 一 / S 二”；tokens 可读“token 数”，第一次可解释为模型输入量。
\item S0 原始自动裁判为 16/30，应用既有证据裁决后为 15/30；S2 两者均为 19/30。正文统一展示裁决后结果。
\item 60 条输入指纹核对与评分离线复算已经完成；总体和分类分数与汇总一致。
\item 分工表展示目前大致贡献，成员沿用 GitHub 账号。
\item 协作方式统一表述为 GitHub 协作、微信群沟通和线下讨论，结合共同题集与实验链路说明。
\end{itemize}

\newpage
\talk{1}{题目与汇报范围}{15 秒\quad 0:00--0:15}{\speechOne}

\vspace{7mm}
\talk{2}{研究问题}{35 秒\quad 0:15--0:50}{\speechTwo}

\newpage
\talk{3}{六篇论文的系统性提炼}{60 秒\quad 0:50--1:50}{\speechSynthesis}

\vspace{7mm}
\talk{4}{StructMem 方法}{40 秒\quad 1:50--2:30}{\speechThree}

\newpage
\talk{5}{实验设计}{40 秒\quad 2:30--3:10}{\speechFour}

\vspace{7mm}
\talk{6}{实际进展}{30 秒\quad 3:10--3:40}{\speechFive}

\newpage
\talk{7}{初步结果与输入成本}{55 秒\quad 3:40--4:35}{\speechSix}

\vspace{7mm}
\talk{8}{共同案例与论文视角}{45 秒\quad 4:35--5:20}{\speechSeven}

\newpage
\talk{9}{创新性论证与后续方案}{45 秒\quad 5:20--6:05}{\speechEight}

\vspace{7mm}
\talk{10}{六人分工与贡献}{40 秒\quad 6:05--6:45}{\speechNine}

\newpage
\talk{11}{小组协作方式}{50 秒\quad 6:45--7:35}{\speechCollaboration}

\vspace{7mm}
\talk{12}{阶段结论与下一步}{25 秒\quad 7:35--8:00}{\speechTen}
\end{document}
